% =====================================================================
%  CHAPTER 1C: 碳水化合物 — 参考：0310营养参考_extracted.txt (PPT原文)
% =====================================================================
\chapter{碳水化合物}
\chapterintro{掌握碳水化合物分类（单/双/寡/多糖）和膳食纤维（可溶/不可溶）的定义与功能；掌握碳水化合物的五大生理功能（供能/抗生酮/节约蛋白质/构成组织/解毒）；掌握血糖指数（GI）和血糖负荷（GL）的定义、计算和分类；熟悉参考摄入量与食物来源。}

\section{碳水化合物的分类与膳食纤维}

\begin{keybox}
\textbf{碳水化合物的分类：}
\begin{itemize}[leftmargin=*]
    \item \imp{单糖（Monosaccharide）}：葡萄糖（glucose）、果糖（fructose）、半乳糖（galactose）、糖醇类
    \item \imp{双糖（Disaccharide）}：蔗糖（sucrose）、麦芽糖（maltose）、乳糖（lactose）、海藻糖（trehalose）
    \item \imp{寡糖（Oligosaccharide）}：棉子糖（raffinose）、水苏糖（stachyose）
    \item \imp{多糖（Polysaccharide）}：可消化（淀粉starch/糊精dextrins/糖原glycogen）和不可消化（抗性淀粉/膳食纤维）
\end{itemize}

\textbf{膳食纤维（Dietary Fiber）：}食物中不能被消化利用的多糖和木质素。分类（按水溶性）：
\begin{itemize}[leftmargin=*]
    \item \imp{不溶性纤维}：纤维素、半纤维素、木质素
    \item \imp{可溶性纤维}：果胶、树胶、一部分半纤维素
\end{itemize}

\textbf{膳食纤维的生理功能：}减少便秘/预防大肠疾病；预防癌症；预防心血管疾病和胆石症；防治糖尿病；防治肥胖。
\end{keybox}

\section{碳水化合物的生理功能}

\begin{defbox}
\textbf{1. 提供能量}

\textbf{2. 构成身体组织}

\textbf{3. 辅助脂肪氧化——抗生酮作用：}脂肪酸在体内分解代谢时产生的乙酰基需与碳水化合物代谢产生的草酰乙酸结合才能进入三羧酸循环而最终被彻底氧化。当碳水化合物不足时，因草酰乙酸不足使得脂肪酸不能被彻底氧化分解而产生过多酮体，会发生酮症酸中毒。反之，当碳水化合物充足时可防止酮症酸中毒的发生，这种作用称为\imp{抗生酮作用}。

\textbf{4. 节约蛋白质作用：}如碳水化合物充足，人体首先利用碳水化合物供给能量，从而减少了蛋白质单纯作为供给能量的分解代谢，有利于改善和提高蛋白质的利用，并增强人体内氮的贮存量。这就是所谓的\imp{节约蛋白质作用}。

\textbf{5. 帮助肝脏解毒}
\end{defbox}

\section{血糖指数（GI）与血糖负荷（GL）}

\begin{keybox}
\textbf{血糖指数（Glycemic Index, GI）：}含50g碳水化合物的食品引起的血糖反应曲线下的面积与含等量碳水化合物的标准食品血糖反应之比。以葡萄糖或白面包为标准。

$$\text{GI} = \frac{\text{摄入含50g CHO食物后2h血糖曲线下面积}}{\text{摄入50g葡萄糖(或白面包)后2h血糖曲线下面积}}$$
\imp{低GI食物：GI $\leq$ 55}；中GI：55 $<$ GI $\leq$ 70；\imp{高GI食物：GI $>$ 70}

\textbf{血糖负荷（Glycemic Load, GL）} = GI × 可利用碳水化合物（CHO $-$ 膳食纤维）。GL同时考虑了食物的升糖速度（GI）和实际摄入的可利用碳水化合物量，能更准确反映一餐/一份食物对血糖的真实影响。
\end{keybox}

\section{参考摄入量与食物来源}

\begin{defbox}
成人CHO：EAR为120g/d，\imp{AMDR占总能量的50\%$\sim$65\%}。添加糖限制：建议不超过50g/d，最好低于25g/d。酒精摄入量不超过15g/d，任何形式的酒精对人体健康都无益处。

食物来源：粮谷类及制品（核心主食来源）、块根(茎)类蔬菜、杂豆类及制品、水果及制品、纯糖类及其制品。
\end{defbox}

\section{复习题}
\begin{enumerate}[leftmargin=*, itemsep=8pt]
    \item \textbf{【名词解释】}膳食纤维；血糖指数（GI）；血糖负荷（GL）；抗生酮作用；节约蛋白质作用
    \item \textbf{【简答题】}简述碳水化合物的分类及膳食纤维的生理功能。
    \item \textbf{【简答题】}试述碳水化合物的抗生酮作用和节约蛋白质作用的机制。
    \item \textbf{【简答题】}简述GI的定义、分类标准和GL的计算公式及意义。
\end{enumerate}

\subsection*{参考答案要点}
\begin{enumerate}[leftmargin=*, itemsep=5pt]
    \item 见正文。\textbf{膳食纤维}：不能被消化利用的多糖和木质素，分不溶性（纤维素/半纤维素/木质素）和可溶性（果胶/树胶）。\textbf{GI}：含50g CHO食物血糖曲线面积/标准食品血糖曲线面积。\textbf{GL}=GI×可利用CHO。
    \item 单糖/双糖/寡糖/多糖（可消化：淀粉/糊精/糖原；不可消化：抗性淀粉/膳食纤维）。膳食纤维功能：防便秘/防癌/防心血管病/防糖尿病/防肥胖。
    \item 抗生酮：脂肪酸分解需草酰乙酸→三羧酸循环，CHO不足→草酰乙酸不足→酮体过多→酮症酸中毒。节约蛋白质：CHO充足→优先供能→减少蛋白质分解供能→改善氮平衡。
    \item GI定义见第1题。低GI≤55，中55-70，高>70。GL=GI×可利用CHO(CHO-膳食纤维)，兼顾升糖速度与实际CHO摄入量。
\end{enumerate}
\newpage
